% Written by travel_notes 0.2.0 from a Markdown note.
\documentclass{amsart}
\usepackage[T1]{fontenc}
\usepackage{amssymb}
\usepackage{hyperref}

\newtheorem{theorem}{Theorem}
\newtheorem{lemma}[theorem]{Lemma}
\newtheorem{proposition}[theorem]{Proposition}
\newtheorem{corollary}[theorem]{Corollary}
\newtheorem{conjecture}[theorem]{Conjecture}
\newtheorem{claim}[theorem]{Claim}
\theoremstyle{definition}
\newtheorem{definition}[theorem]{Definition}
\newtheorem{example}[theorem]{Example}
\newtheorem{exercise}[theorem]{Exercise}
\theoremstyle{remark}
\newtheorem{remark}[theorem]{Remark}

\begin{document}

\title{Math next to Markdown}
\maketitle

Inline \(a_1 + a_2 = b_{12}\) and \(c*d*e\) and \(f\_g\): underscores, asterisks and
backslashes inside math are left alone.

Emphasis around math: \emph{the value \(x^2\) is large} and \textbf{\(y\) is small}.

Display math on its own lines:

\[
\int_0^1 x^2 \, dx = \frac{1}{3}
\]

An aligned environment with a label:

\begin{align}
a &= b + c \label{eq:first} \\
d &< e \& f
\end{align}

A macro: \newcommand{\R}{\mathbb{R}} \(f\colon \R \to \R\).

Dollar signs in prose: it costs \$5, then \(10, or 15\) in total. An escaped \$ sign.

\begin{itemize}
\item a list item with \(\alpha_i\)\begin{itemize}
\item a nested item with
\[
\sum_{k=1}^n k = \frac{n(n+1)}{2}
\]

\end{itemize}


\end{itemize}

\begin{quote}
A quotation with \(\ell\)-isogenies
and
\[
\Phi_\ell(X, Y) = 0
\]
on the next line.

\end{quote}

\begin{center}
\begin{tabular}{rr}
\hline
\(x\) & \(x^2\) \\
\hline
2 & \(4\) \\
\hline
\end{tabular}
\end{center}

Code keeps its dollars: \texttt{\$not math\$}.

\begin{verbatim}
$$ also not math $$
\end{verbatim}

Unterminated: \$x + y and \$\$z.

\end{document}
